% ════════════════════════════════════════════════════════════════════════════
\section{Classification of existing designs}
\label{sec:classify}
% ════════════════════════════════════════════════════════════════════════════

Corollary~\ref{cor:reach} gives the classification axis: a design that clears
Criterion~\ref{prin:nga} does so either by \emph{removing the race} or by
\emph{removing $\Attacker$'s reach over the racer}. This section places the
shipping designs on it --- stating each verdict and its reason, with the full
analyses in Appendix~\ref{app:classify}.

\subsection{Time-lock-rescue vaults, and the reachable-racer failure}

\textbf{Rewind Bitcoin}~\cite{rewind-bitcoin} holds funds in a time-lock vault:
unlocking starts a multi-day countdown within which the owner \emph{or a trusted
delegate} may press ``Rescue'' and sweep to a cold panic address. Against
\emph{remote} theft this is sound. But the spend capability sits behind \emph{only
a delay}, keys on the seized device and no recall floor, so a wrench attacker
holds a feasible path --- unlock, wait out the countdown. That is a race, and
Rewind's answer is exactly \ref{r:remotewinner}: a remote delegate the attacker
cannot reach presses Rescue. A legitimate clearance whose cost is that
\textbf{custody is no longer individual}, and which holds only while the delegate
is genuinely remote (Remark~\ref{rem:escape-closes}).

\textbf{BitVault}~\cite{bitvault} is the same species. No authoritative
specification is published; the account rests on the vendor's pages, a
self-attributed concept whitepaper~\cite{bssl} and correspondence, which do not
agree in every particular. We note this to bound the claim, not as an allegation:
we have no evidence that any material was withdrawn. Taking that whitepaper as
valid, the spending paths carry timelocks --- a primary needing a custodian
co-signature, a one-year fallback on the user's own two keys, a three-year
custodian path --- and both cases fail. Before the one-year date $\Attacker$
coerces $A$ and $B$ and kills, the only party who would tell the gatekeeper to
withhold being $\Victim$: removal is what buys settlement (Lemma~\ref{lem:drl}).
After it he coerces $A$ and $B$ and spends at once, ungatekept --- no delay, no
notification, nothing to cancel. \S\ref{app:classify-bitvault} gives the sources
and their limits. A caveat sharpens \ref{r:remotewinner} itself: the
winning racer must be not merely remote but \emph{un-coercible} --- an autonomous
watchtower, not a human friend, since coercion routes through the victim to a
person and a humane delegate pays (\S\ref{sub:empirical}).

\subsection{The industry recommendation set}
\label{sub:industry}

The same structure appears in the \emph{standard advice}. Threat-intelligence
reporting states the objective in terms close to ours --- custody should ``remain
resilient even when a human is under duress'' --- and recommends withdrawal
delays, transaction-size thresholds, staged vault architecture, emergency freeze
processes, and multisig with geographically distributed
signers~\cite{certik-intel3d-h1-2026}. Criterion~\ref{prin:nga} does not reject
this advice; it \emph{partitions} it, and that partition is the framework's
practical content. Each control interposes a delay or a veto, hence defines a
competing racer, so the only question is whether $\Attacker$ can reach whoever
holds the release. Where that authority rests with a remote, institutionally
separated party, the control clears the bar by \ref{r:remotewinner} at that
route's standing cost; where the delay is one the \emph{holder} can release, or
the freeze authority is a device or person on the scene, it leaves a
feasible-but-slower path --- the gray outcome the Criterion forbids, and the
configuration in which the control is worse than its absence. It is the
reachability question, not the length of the delay, that separates a delay which
buys rescue from one which buys $\Attacker$ a reason to eliminate the other
racer.

\subsection{Self-destruct (wallet-erase) PINs}
\label{sub:selfdestruct}

The designs above interpose a delay or a veto. A large and widely shipped class
does something categorically different: it \emph{erases}. Blockstream's Jade
documents such a PIN under both names its help center uses --- \emph{duress} and
\emph{wallet-erase} --- as an alternate PIN that ``automatically delete[s] the
encrypted wallet information from Jade if entered,'' ``especially useful in the
event of a physical threat, as you can provide this PIN to an attacker''; the
device then reports only ``Internal Error,'' so the erasure reads as a
malfunction~\cite{jade-duress-pin,jade-erase-pin}. Grant it the strength the delay architectures lack: it is
\emph{instantaneous and covert}, so $\Attacker$ cannot intervene. The claim nonetheless decomposes into a real half and an obscurity bet
--- \textbf{(a)} the \emph{erasure}, which happens whatever $\Attacker$ knows, and
\textbf{(b)} the \emph{belief} that nothing further remains, manufactured by the
``Internal Error'' report. Only the erasure is Kerckhoffs-independent, and how much coercion
resistance it carries alone is answered by the vendor's own instructions:
\emph{none}. The same documentation warns that afterwards ``there is no way to
regain access to your funds except by restoring using your recovery phrase,'' and
directs the holder to ``make sure you have a backup before
proceeding''~\cite{jade-erase-pin}. The
feature \emph{presupposes} a surviving backup: a holder who follows the
documentation has one, and one who does not has destroyed the stash rather than
defended it. What erasure removes is therefore not the secret but the
\emph{device path} to it --- a path that cost $\Attacker$ nothing to lose, since
$\Victim$ is left alive, present and still able to restore, and one contributing
nothing to the floor (Remark~\ref{rem:shallow-null}). The route that remains
runs through $\Victim$'s memory of a seed phrase, or of where a backup is kept ---
\textbf{the erasure has deleted the path that required hurting no one and retained
the path that does.}

That also makes it the worse of the two obscure bets it is sold beside, by an
ignorance count (Remark~\ref{rem:two-ignorances}): a decoy needs $\Attacker$
ignorant only of decoys, and granted that it \emph{terminates the encounter},
whereas a self-destruct needs him ignorant of the feature \emph{and} of NNPP's
consequence --- so even total success at the first leaves the encounter open,
``this device is broken'' inviting \emph{where is your backup?} from a $\Victim$
still present and still knowing. It is the one obscure scheme with no branch on
which it wins, and its own vendors document the regress: Coldcard's firmware
guidance, for the case where the duress wallet is disbelieved, advises escalating
to the \emph{brickme} PIN~\cite{coldcard-trick-pins,coldcard-pin-entry} --- an
irreversible erasure performed with $\Victim$ still in the room. The class lands
in the \textbf{naive} tier. Appendix~\ref{app:classify} records what each of the
three vendors shipping the pair actually \emph{claims} for it (two decline the
coercion claim on one of their two mechanisms), the two bluffs an informed
$\Attacker$ can play instead of entering any PIN, the companion decoy --- sold as
coercion protection and coached for believability
(Remark~\ref{rem:decoy-market}) --- and the \emph{Obscurity Paradox}.

\subsection{TGPO, and the honest comparison}

TGPO clears by the \emph{other} route: the secret is tacit and never on the
device, so a seizure yields nothing feasible and there is no racer to reach. Read
through the Four Properties the comparison is a covering-set statement --- each
shipping design trades at least one property TGPO keeps, and one trades a property
and \emph{still} fails.
\begin{center}\footnotesize
\renewcommand{\arraystretch}{1.3}
\begin{tabularx}{\linewidth}{@{}l cccc X@{}}
\toprule
\textbf{Design} & \textbf{(\ref{p:know})\,Know.} & \textbf{(\ref{p:indiv})\,Indiv.} & \textbf{(\ref{p:nonob})\,Non-ob.} & \textbf{(\ref{p:coerc})\,Coerc.} & \textbf{What it trades} \\
\midrule
\textbf{TGPO} (this work) & \cmark & \cmark & \cmark & \cmark & holds all four (tacit; \ref{r:removerace}) \\
Rewind Bitcoin               & \xmark & \xmark & \cmark & \cmark & individual custody: a remote delegate (\ref{r:remotewinner}) \\
BitVault                     & \qmark & \xmark & \qmark & \xmark & individual custody both ways --- denial blocks the user, collusion spends past three years \\
Casa                         & \cmark & \xmark & \cmark & \cmark & custody shared with remote co-signer\\
Self-destruct PIN            & \xmark & \cmark & \xmark & \xmark & non-obscurity: the erasure must read as malfunction (Coldcard, Jade) \\
\bottomrule
\end{tabularx}
\end{center}
\noindent Row~(\ref{p:know}) is the covering point: TGPO is the \emph{only}
knowledge-based row, every competitor keeping a \emph{transmissible} secret on a
device. Casa's collaborative multisig trades individual custody, a
coercion-resistant delegate being remote \emph{by definition}~\cite{casa}. The axis tilts further to \ref{r:removerace}
for a blunter reason: the delegated rows (Rewind, Casa) are not self-custody at all
(Remark~\ref{rem:r2-cost}).

\subsection{Empirical grounding in the registry}
\label{sub:empirical}

The two stipulations leaning on the physical-attack registry~\cite{lopp-attacks}
--- the coercion-time budget and the murder-cost weighing of
\S\ref{sec:residual-dr} --- admit directional support. The registry records only
\emph{date, victim, location, headline}, so we code its $352$ incident headlines
(at commit \texttt{9a4a62a}) by keyword; the coding is coarse and a reproducible script
accompanies the paper.\footnote{\texttt{registry-analysis/} gives the keyword
patterns, per-year counts and caveats; Appendix~\ref{app:classify} gives the
method, the coding tables, the jurisdictional confounds and the case analyses.}
Four findings bear on the paper.

\emph{The attack is frequent and worsening}: recorded incidents climb from a
handful a year before 2017 to a peak of $85$ in 2025, and a separately sourced
count --- threat-intelligence reporting from law-enforcement disclosures and court
documents --- agrees, with $72$ verified incidents in 2025, up $75\%$ year on
year~\cite{certik-skynet-2026}, and $52$ in the first half of 2026 against $39$ a
year earlier~\cite{certik-intel3d-h1-2026}. Reported \emph{home invasions} there
rose from $1$ to $20$ of $52$: coercion is moving to where the holder lives, the
setting the model assumes.

\emph{The terminal move is real}: at least $5.1\%$ ($18/352$) of headlines record
a \emph{killed} or \emph{murdered} victim. A homicide is usually reported as a
homicide rather than as a ``Bitcoin attack,'' so this is a \emph{lower} bound ---
yet even so the outcome Lemma~\ref{lem:drl} predicts is attested in practice.
That censoring compounds with a second: obscurity does not advertise itself, so no
victim's exposure to an obscurity-dependent design is on record either. The
numerator is censored by how the terminal event is classified and the covariate by
what vendors do not say, which is why the pattern cannot be read off the data and
the bar has to be derived instead (Remark~\ref{rem:anosognosia}).

\emph{The mix is not universal, and the Lemma's bite tracks it}: by the ratio of
long-hold to short-hold codings, France and the United States --- the two largest
national subsets, of near-identical size ($61$ and $59$) --- have \emph{inverted}
compositions, $7.0$ against $0.7$. Media inflation explains a \emph{level} but not
this \emph{ordering}, and the difference matters for scope rather than validity: a
hand-over-now robbery approximates the clean $p_A = 1$ endpoint, a multi-hour hold
ending in deliberate release is exactly the gray outcome. The Criterion does its
work where long holds dominate --- and France alone accounts for a majority of the
incidents recorded in 2026.

\emph{The attacker often does not need the holder at all}: in $12$ of the $352$
incidents ($3.4\%$, hand-coded and audited) the person seized was \emph{not} a
holder but a relative held to compel one, eleven in 2023 or later and nine in
France. This is \emph{not} the hostage token of \S\ref{sub:hostage} --- the
leverage is a person rather than a seizable object, so devicelessness does not
reach it --- but clause~\ref{d:nodelegatecoerce} failing in the field, and it
bounds \ref{r:remotewinner}: the delegate a holder is most likely to appoint is
precisely the party being seized, and a reachable delegate is not a second racer
but a second victim.

Two fielded cases sharpen this, both analysed in Appendix~\ref{app:classify}: the
Balland abduction~\cite{balland}, whose rescue and recovery ran through exactly
the state and centralized channels clause~\ref{d:nodelegatecoerce} assumes away;
and the RATIRL torture~\cite{ratirl}, where a denial was \emph{pre-discredited}
because denial is the norm, so the rational continuation was \emph{more} coercion.
The second fixes the \emph{sign}: obscurity does not merely fail to protect, it
can convert a robbery into a torture.
